\section{Which Neumann problems have a gauge?}
For $\mathcal L u=-\nabla\cdot(D\nabla u)+Cu$ and a constant $a$,
\[
  \mathcal L(u+a)=\mathcal L u+Ca.
\]
Thus Neumann boundary conditions alone do not give a gauge.  If
$C\not\equiv0$, shifting the solution changes the equation.  Subtracting
the mean of the error would hide a real error; for example, the constant
field $u_h=7$ compared with $u=0$ has error $7$, not zero.  For
$C\equiv0$ and all-Neumann data, error is compared modulo constants using
the visible-cell volume weights.  This distinction is independent of
source quadrature.  The regression \code{testEbNeumannSemantics} checks
both cases, including \code{SourceQuad=3} and the well-conditioned scaled
problem $D=10^{-18}$, $C=f=10^{-14}$, $g=0$, whose solution is $u=1$.

\section{Compatibility after every refinement}
Let $V_i=h_{x,i}h_{y,i}$ denote a visible cell's area and $A_F$ a
physical boundary face's length.  With outward derivative data $g_F$,
the singular conservative equations require
\[
 S+J=0,\qquad S=\sum_{i\in\vis}V_i f_i^{\rm raw},\qquad
 J=\sum_{F\subset\partial\Omega} A_F D_F g_F.
\]
The implementation computes
\[
 c_H=\frac{S+J}{\sum_{i\in\vis}V_i},\qquad
 f_i^{\rm projected}=f_i^{\rm raw}-c_H.
\]
It records $c_H$ in \code{H.compatibility.sourceShift} and $S+J$ in
\code{H.compatibility.unprojectedDefect}.  Each level's
\code{sourceShift} records the constant already removed from its
\code{FTrue}.  A fresh level starts at zero.  Reprojection restores the
old shift in the composite sum and applies the difference between the
new and old shifts, so every level represents the same projected source.
Without this bookkeeping, successive refinements introduce different
constant source perturbations on different levels even when the final
global mass defect is small.

For smoothly compatible continuous data, midpoint source and boundary
quadrature normally make $c_H=O(h_{\max}^2)$ on a controlled mesh family.
For incompatible data it may be $O(1)$: projection then solves a changed
problem.  Tensor Gauss quadrature improves source integration but is not
exact for arbitrary functions and does not improve the separately
sampled Neumann boundary flux automatically.

\section{A multidimensional QIF sign counterexample}
Take constant $D>0$ and a west coarse--fine face away from tangential
end exceptions.  In the padded fine array, \code{ebInterpCf} writes
\[
 Q_{1,3}=\frac8{15}Q^c_{1,2}+\frac23Q_{2,4}-\frac15Q_{3,5}.
\]
The operator row for interior value $Q_{2,3}$ contains the term
$-D Q_{1,3}/h_x^2$.  Therefore the coefficient of the distinct interior
unknown $Q_{3,5}$ in that row is
\[
  +\frac{D}{5h_x^2}>0.
\]
That unknown is not one of the row's four ordinary nearest neighbors,
so no nearest-neighbor diffusion coefficient cancels it.  This remains
true for constant coefficients and square cells.  Consequently the
fully ghost-eliminated multidimensional QIF matrix is not, in general,
a $Z$-matrix.  A one-dimensional normal-stencil sign argument cannot be
applied to this diagonal stencil.  Any comparison principle, inverse
positivity result, or energy estimate must be justified for the actual
composite operator; flux conservation alone does not establish it.

The existing \code{convergence\_theory.tex/pdf} is retained separately.
Its multidimensional sign claim and the theorems relying on that claim
require further proof work; they should not be treated as verified
theorems for this implementation.  The measured convergence tables are
finite numerical evidence.  They neither prove a uniform V-cycle bound
nor control arbitrary interface geometry, coefficient contrast, or an
unbounded number of patches.

\section{What the diagnostics establish}
The composite residual includes the flux correction and measures
algebraic error.  At any synchronized state, the conservative mass
defect equals minus the volume sum of that residual, up to roundoff.
It can therefore be small through cancellation while the maximum
residual remains large.  Always report both.  The error norm adds an
independent comparison with a correctly manufactured solution.

The \code{ulap} indicator is
$\sqrt{(\delta_x^2u)^2+(\delta_y^2u)^2}$, with undivided differences;
it is a refinement heuristic, not a rigorous error bound.  Proper
nesting and minimum-width handling may drop tags.  Moreover, the driver
only appends a new finest level: it does not revisit dropped tags on
older levels or coarsen an existing hierarchy.  Inspect the resulting
visible-cell coverage when a particular region must be resolved.

\section{Reproducible command-line checks}
From the repository root on the audited macOS installation:
\begin{lstlisting}[language=bash]
/Applications/MATLAB_R2026a.app/bin/matlab -batch "maxNumCompThreads(2); cd('EllipticBSAM/tests'); unitKernels; runAllTests; auditExtras; stressInterfaces; testNeumannCompat; testEbNeumannSemantics"
\end{lstlisting}
Use a separate MATLAB path for each solver's test directory: both 2D
packages define \code{mmsCases}.  The September 2026 evidence under
\code{Audit/2026-09-10/} distinguishes existing suite results from the
new regression tests.  Historical all-Neumann, positive-reaction error
tables may have removed the error mean; rerun them before comparing
absolute errors with the corrected diagnostics.
